\subsection{Threat Model and Operational Setting}
\label{subsec:threat-model}

We consider a contested Cyber-EMS environment in which a defender operates a sensor (a software-defined radio with antenna array) that observes a shared RF band. An adversary intermittently emits RF waveforms intended to enable a downstream electronic or cyber effect. For example, a burst pattern that probes a receiver's automatic gain control, a sustained emission used as a denial-of-service primitive, a hopping pattern used to evade fixed-channel monitoring, or a replay segment used to inject a previously captured cooperative signal. We refer to these emissions as \emph{preliminary actions} (PAs) following~\cite{trott2026model}, and we treat the \emph{full} attack which is the operational effect that the PA enables as occurring later in time.

Two assumptions are central. First, the adversary's PA repertoire is \emph{partially} known: the defender has labeled training examples for a finite set of PA classes, but the adversary may at any time deploy a PA whose waveform-level behavior falls outside that set. Second, the defender's response is \emph{costly} and \emph{time-bounded}: each candidate response (continued monitoring, additional sensing, deception, escalation) consumes resources and is only valuable if executed before the adversary's full tactic completes. The objective is therefore not classification accuracy in isolation, but the joint quantity \emph{probability of correctly characterizing the PA}~$\times$~\emph{duration of the resulting QR-CWoS}.

\subsection{Preliminary Action Taxonomy}
\label{subsec:pa-taxonomy}

Following the pilot taxonomy of~\cite{trott2026model}, we define four behaviorally distinct PA classes that serve as the trained label space:

\begin{itemize}
    \item \textbf{PA2: Burst.} The emitter transmits in many short on-intervals separated by quiet pauses; the defining feature is the temporal pattern of burst lengths and inter-burst gaps rather than the carrier itself.
    \item \textbf{PA3: Sustain.} The emitter transmits nearly continuously across the observation window with very few start/stop transitions, presenting as a steady occupied channel rather than discrete events.
    \item \textbf{PA4: Hop.} The emitter transmits while switching among a small set of frequencies, dwelling briefly on each and often revisiting previously used frequencies within the same window.
    \item \textbf{PA8: Replay.} The emitter records a short segment of a transmission and replays that exact segment repeatedly; replays are bit-for-bit copies of a waveform segment separated by short quiet gaps, which distinguishes them from newly generated bursts.
\end{itemize}

In addition, we hold out at least one PA class entirely from training to serve as the \emph{unseen-tactic} condition. The held-out PA is present in the over-the-air (OTA) test stream but never in training; the pipeline must label it \texttt{unknown} rather than collapse it into the nearest of \{PA2, PA3, PA4, PA8\}.

\subsection{Signal Acquisition and Preprocessing}
\label{subsec:preprocessing}

Each captured emission is recorded as a complex baseband I/Q stream and segmented into fixed-length analysis windows. For each window we compute (i)~a short-time Fourier spectrogram to expose the time-frequency structure that distinguishes burst from sustain and hop from replay, and (ii)~an event-detection trace that thresholds the per-bin energy to recover the on/off and dwell statistics on which PA2, PA3, and PA4 are partially defined. Windows are validated by the OTA validation pipeline of~\cite{trott2026model} to ensure TX/RX synchronization and to reject windows whose payload integrity has been compromised by channel effects; only validated windows enter the classifier. The output of preprocessing is a tensor representation $x \in \mathbb{R}^{C \times T \times F}$ together with a small vector of hand-engineered event-statistics features used as auxiliary input.

\subsection{Open-Set PA Classification with varMax}
\label{subsec:varmax}

The PA classifier is a deep neural network $f_\theta$ producing logits $z = f_\theta(x) \in \mathbb{R}^K$ over the $K=4$ trained PA classes. Rather than reading off $\arg\max \mathrm{softmax}(z)$, which is well known to be overconfident on out-of-distribution inputs, we apply the varMax open-set decision rule~\cite{baye2024varmax,broggi2025varmax}, which uses the \emph{variance} and the \emph{maximum} of the logit vector jointly as the basis for a confidence-aware decision. Concretely, at inference time we compute
\begin{equation}
\hat{y} =
\begin{cases}
\arg\max_k z_k & \text{if } g(z) \geq \tau \\
\texttt{unknown} & \text{otherwise}
\end{cases}
\label{eq:varmax-rule}
\end{equation}
where $g(z)$ is the varMax confidence statistic that combines the spread of the logits with their peak value, and $\tau$ is a calibrated threshold chosen on a held-out validation split that contains both in-distribution PAs and synthetic perturbations. This produces a five-way output: PA2, PA3, PA4, PA8, or \texttt{unknown}. We also retain the continuous confidence value $g(z)$ itself, which we expose to the downstream policy so that ``low-confidence PA2'' and ``high-confidence PA2'' induce distinct policy states.

For the application reported in~\cite{trott2026model}, we retain the network architecture validated on the pilot OTA dataset and adapt only the final decision layer. This isolates the effect of the open-set rule from architectural choices and lets us attribute downstream gains to the novelty-aware decision boundary rather than to representation changes.

\subsection{DQN-based Response Selection}
\label{subsec:dqn}

The varMax output is a label and a confidence, not an action. To map outputs to actions we adopt the DQN-IDS formulation of~\cite{tiwaridqn}, originally developed for open-set network intrusion detection, and adapt it to the EW response setting.

\paragraph{State.} At each decision step $t$, the agent observes
\begin{equation}
s_t = \bigl(\hat{y}_t,\ g(z_t),\ \phi_t,\ h_{t-1}\bigr),
\label{eq:state}
\end{equation}
where $\hat{y}_t \in \{\text{PA2, PA3, PA4, PA8, \texttt{unknown}}\}$ is the varMax verdict, $g(z_t)$ is its confidence, $\phi_t$ is the auxiliary event-statistics vector from preprocessing, and $h_{t-1}$ is a short history summary of recent verdicts (rolling counts and dwell time per class) so the policy can recognize behavioral patterns that span multiple windows. Note that DQN can functionally be an alternative to be used in the place of varMax. For such cases, we will compare the two algorithms isolated later in this paper.

\paragraph{Action space.} The agent selects an action $a_t \in \mathcal{A}$ from a discrete set of operationally meaningful responses. In our deployment $\mathcal{A}$ contains: continue passive observation; request a higher-rate capture on the same band; allocate the SDR's directional beam to localize the emitter; deploy an EMS deception waveform consistent with the framework of the parent project; and escalate to operator review. The set is small enough for tabular Q-learning to be infeasible only because of the continuous confidence component, which is why a function-approximation DQN is appropriate.

\paragraph{Reward.} The reward function is the place where QR-CWoS enters training. Rather than rewarding raw classification accuracy, we reward the agent for \emph{opening and maintaining} a usable window. For each episode the reward is the sum of (i)~a positive term proportional to the duration between PA detection and adversary tactic completion (the open window), (ii)~a negative term penalizing actions that unnecessarily consume sensing or deception resources, and (iii)~a large negative terminal reward if the adversary's tactic completes without the defender having selected a response other than passive observation. The \texttt{unknown} verdict is \emph{not} directly penalized; what is penalized is reaching the end of the episode having ignored an unknown.

\paragraph{Training.} We train the DQN with experience replay and a target network, following~\cite{tiwaridqn}. Episodes are constructed by drawing trajectories from the OTA dataset and interleaving held-out PAs at controlled rates, so the policy experiences both familiar and unseen tactics during training. To prevent the policy from overfitting to a particular adversary profile, we vary the inter-PA timing distribution and the prevalence of each PA across episodes.

\subsection{QR-CWoS Quantification}
\label{subsec:qrcwos}

The pipeline outputs a verdict and an action at every analysis window, but the operational quantity of interest is the QR-CWoS itself. We define the window for an episode as the interval $[t_d, t_c]$, where $t_d$ is the time at which the pipeline first emits a non-passive response with confidence above a deployment threshold, and $t_c$ is the time at which the adversary's full tactic would have completed in the absence of intervention (this completion time is recoverable in our OTA experiments because we control the emitter). Three derived quantities are reported:

\begin{enumerate}
    \item \textbf{Window duration}, $\Delta t = t_c - t_d$, the time during which the defender has acted but the adversary has not yet completed, i.e., the raw width of the QR-CWoS.
    \item \textbf{Window reliability}, the fraction of episodes for which $\Delta t$ exceeds a stated operational minimum, conditioned on whether the PA was seen or unseen at training time. This is the resilience axis of QR-CWoS: a wider mean window is not useful if the variance is high.
    \item \textbf{Window quality under uncertainty}, the conditional distribution of $\Delta t$ given the varMax confidence $g(z_t)$ at $t_d$. This is the quantification axis: it tells the operator how much window to expect for a given confidence level, and it lets the DQN's value function be interpreted as an estimate of expected window value.
\end{enumerate}

Reporting all three quantities rather than only mean accuracy is what distinguishes a QR-CWoS evaluation from a conventional classification evaluation. It is also what makes the pipeline auditable: an operator can ask, for any deployment, ``what window does this system actually deliver against tactics it has not seen, and how reliably?'', and the pipeline will answer in seconds rather than in classes.